\documentclass[12pt,a4paper,UTF8]{ctexart}
\usepackage[a4paper,top=22mm,bottom=22mm,left=22mm,right=22mm,headheight=15pt]{geometry}
\usepackage{amsmath,amssymb,mathtools,bm}
\usepackage{graphicx,xcolor,booktabs,longtable,array,enumitem,fancyhdr,caption,float,microtype,setspace}
\usepackage[hidelinks,unicode]{hyperref}
\usepackage[most]{tcolorbox}
\definecolor{ContestBlue}{HTML}{17365D}
\hypersetup{pdftitle={第086套},pdfauthor={}}
\setstretch{1.24}
\setlength{\parindent}{2em}
\setlength{\parskip}{0.35em}
\setlength{\emergencystretch}{2em}
\setlist{itemsep=0.35em,topsep=0.4em}
\captionsetup{font=small,labelfont=bf,justification=centering}
\setcounter{secnumdepth}{0}
\pagestyle{fancy}
\fancyhf{}
\fancyhead[L]{\small 第086套}
\fancyfoot[C]{\thepage}
\renewcommand{\headrulewidth}{0.4pt}
\ctexset{
  section={format=\Large\bfseries\color{ContestBlue},beforeskip=1.2em,afterskip=0.7em},
  subsection={format=\large\bfseries\color{ContestBlue},beforeskip=1.0em,afterskip=0.5em},
  subsubsection={format=\normalsize\bfseries,beforeskip=0.8em,afterskip=0.35em}
}
\makeatletter
\newcommand{\tightlist}{\setlength{\itemsep}{0pt}\setlength{\parskip}{0pt}}
\newcommand{\maxwidth}{\ifdim\Gin@nat@width>\linewidth\linewidth\else\Gin@nat@width\fi}
\newcommand{\maxheight}{\ifdim\Gin@nat@height>0.72\textheight 0.72\textheight\else\Gin@nat@height\fi}
\newcommand{\pandocbounded}[1]{\begingroup\setkeys{Gin}{width=\maxwidth,height=\maxheight,keepaspectratio}#1\endgroup}
\makeatother
\begin{document}
\begin{titlepage}
  \centering
  \vspace*{58mm}
  {\Huge\bfseries\color{ContestBlue} 第086套\par}
  \vfill
  \vspace*{28mm}
\end{titlepage}
\section{第 1 题}\label{ux7b2c-1-ux9898}

如图，在一个足够长的光滑斜面上连接一个边长为 \(d\)
的正方形光滑水平平台，斜面的倾角为 \(30^\circ\)。水平平台上有一质量为
\(m\)
的小球。轻绳的一端系着小球，穿过平台中心的光滑小孔，另一端系在斜面上一质量为
\(M=2m\) 的小物块上，物块可视为质点。绳长为 \(1.5d\)，重力加速度为
\(g\)。

\begin{enumerate}
\def\labelenumi{\arabic{enumi}.}
\tightlist
\item
  若小球绕小孔做半径为 \(d/4\) 的匀速圆周运动，求角速度 \(\omega_0\)；
\item
  在上述运动状态下，突然给物块一个沿斜面向下的冲量，使物块获得速度
  \(v'=\sqrt{2gd}\)，求此时小球指向圆心的径向速度；
\item
  此后若物块始终与斜面接触，求物块沿斜面向下运动的最大位移，结果给出数值解；
\item
  重新回到第（1）问的状态。沿小球初速度方向再给小球一个冲量
  \(I\)，使其切向速度瞬间增大。求当 \(I\)
  超过多大时，小球有可能从正方形平台边缘掉下。
\end{enumerate}

\begin{figure}[H]
\centering
\includegraphics[width=0.502\linewidth,height=\textheight,keepaspectratio,alt={图 1-1　水平平台、斜面、轻绳与两物体的空间位置。}]{../figures/question-01.pdf}
\caption*{图 1-1　水平平台、斜面、轻绳与两物体的空间位置。}
\end{figure}

\newpage

\section{第 2 题}\label{ux7b2c-2-ux9898}

某种物质的电导率是非线性的，其电阻率 \(\rho\) 与电场强度 \(E\) 满足

\[
\rho=\rho_0+AE^2,
\]

其中

\[
\rho_0=1.7\times10^7\,\Omega\cdot\mathrm m,
\qquad
A=1.0\times10^{-3}\,\Omega\cdot\mathrm m^3/\mathrm V^2.
\]

用该物质充满平行板电容器的两块极板之间，极板面积 \(S=1\,\mathrm{m^2}\)。

\begin{enumerate}
\def\labelenumi{\arabic{enumi}.}
\tightlist
\item
  求流过电容器的电流最大值 \(I_{\max}\)；
\item
  极板间距
  \(d=1\,\mathrm{cm}\)，在两极板间施加可变电压，求介质内最大焦耳热功率，并画出
  \(P-U\) 关系图；
\item
  保持电压 \(U_1=2.0\times10^3\,\mathrm V\) 不变，改变极板间距
  \(d\)。求最大功率、达到最大功率时的 \(d_1\)，并定性画出 \(P-d\) 图。
\end{enumerate}

\begin{figure}[H]
\centering
\includegraphics[width=0.9\linewidth,height=\textheight,keepaspectratio,alt={图 2-1　非线性导电介质填充的平行板电容器。}]{../figures/question-02.pdf}
\caption*{图 2-1　非线性导电介质填充的平行板电容器。}
\end{figure}

\newpage

\section{第 3 题}\label{ux7b2c-3-ux9898}

有一根质量分布不均匀的绳子，总长为 \(L\)。绳上各点的线密度
\(\lambda(x)\) 与该点到左端点 \(P\)
的距离成正比。把绳放在水平桌面上，使长度为 \(l_0\)
的一段从光滑桌角竖直下垂，桌高大于绳长。由静止释放。绳与桌面间的动摩擦因数为
\(\mu\)，桌角光滑，不考虑绳在运动中飞起。

\begin{enumerate}
\def\labelenumi{\arabic{enumi}.}
\tightlist
\item
  已知 \(\lambda(x)=kx\)。若要使绳释放后不动，至少需要多大的静摩擦因数
  \(\mu'\)；
\item
  若随后摩擦因数降为
  \(\mu<\mu'\)，从初始状态到绳刚离开桌面，摩擦力对绳做的总功是多少；
\item
  从初始状态到绳刚离开桌面，绳的重力势能减少量是多少；
\item
  绳刚离开桌面时的速度是多少。
\end{enumerate}

\begin{figure}[H]
\centering
\includegraphics[width=0.717\linewidth,height=\textheight,keepaspectratio,alt={图 3-1　线密度随位置增大的绳子跨过光滑桌角。}]{../figures/question-03.pdf}
\caption*{图 3-1　线密度随位置增大的绳子跨过光滑桌角。}
\end{figure}

\newpage

\section{第 4 题}\label{ux7b2c-4-ux9898}

两个相同的线轴对称放置在一个对称山坡的两侧，山坡顶角为
\(150^\circ\)。线轴内筒绕有不可伸长的轻绳，轻绳与内筒之间不打滑，内筒半径是外筒半径的一半。绳子中点
\(P\)
从静止开始竖直向上移动。山坡的摩擦因数足够大，线轴始终做无滑动纯滚动。求当
\(P\) 点速度为 \(v\)、两段绳的夹角 \(2\alpha=60^\circ\)
时，两线轴轴线沿水平方向的接近速度 \(u\)。

\begin{figure}[H]
\centering
\includegraphics[width=0.515\linewidth,height=\textheight,keepaspectratio,alt={图 4-1　对称山坡、线轴与上移绳中点。}]{../figures/question-04.pdf}
\caption*{图 4-1　对称山坡、线轴与上移绳中点。}
\end{figure}

\newpage

\section{第 5 题}\label{ux7b2c-5-ux9898}

在自由空间中，半径为 \(R_0\) 的圆内接正方形四个顶点上各有一个质量为
\(m\) 的质点。其中两个带电 \(+q\)，另两个带电
\(-q\)，相邻顶点电荷异号。初始时四个质点以相同速率沿圆的切线顺时针运动。已知运动过程中，任一质点到圆心
\(O\) 的最小距离为 \(R_1<R_0\)；任意时刻四个电荷仍位于以 \(O\)
为中心的正方形四个顶点。忽略重力，只考虑静电作用。

\begin{enumerate}
\def\labelenumi{\arabic{enumi}.}
\tightlist
\item
  每个粒子的运动轨迹是什么曲线；
\item
  求一个粒子从初始位置运动到距圆心 \(R_1\) 处所需的时间；
\item
  在初态沿两个负电荷所在直线方向、距 \(O\) 点 \(l\gg R_0\)
  处，求电场强度精确到首个非零小量。
\end{enumerate}

\begin{figure}[H]
\centering
\includegraphics[width=0.608\linewidth,height=\textheight,keepaspectratio,alt={图 5-1　相邻异号的四电荷正方形与初始切向速度。}]{../figures/question-05.pdf}
\caption*{图 5-1　相邻异号的四电荷正方形与初始切向速度。}
\end{figure}

\newpage

\section{第 6 题}\label{ux7b2c-6-ux9898}

冥王星可视为半径为 \(R\)、质量为 \(M\)
的球体，表面覆盖一层海水，忽略潮汐影响。冥王星绕 \(y\) 轴自转，角速度为
\(\omega\)，万有引力常量为 \(G\)。

\begin{enumerate}
\def\labelenumi{\arabic{enumi}.}
\tightlist
\item
  对海洋表面不在两极的一小质量 \(\Delta m\)，比较其所受万有引力
  \(F_{\text{万}}\) 与支持力 \(N\) 的大小；
\item
  假设海水总量很少且自转不太快，把海洋表面近似为旋转椭球面，其轴截面为椭圆。求椭圆半长轴与半短轴之比；
\item
  求赤道与两极海水深度差 \(\Delta h\)；
\item
  冥王星自转角速度至少多大，赤道处海水才会飞入太空。
\end{enumerate}

\begin{figure}[H]
\centering
\includegraphics[width=0.703\linewidth,height=\textheight,keepaspectratio,alt={图 6-1　冥王星固体表面与旋转海洋等势面。}]{../figures/question-06.pdf}
\caption*{图 6-1　冥王星固体表面与旋转海洋等势面。}
\end{figure}

\newpage

\section{第 7 题}\label{ux7b2c-7-ux9898}

制作一个水火箭。塑料瓶体积为 \(V_0\)，瓶中装有体积 \((1-\alpha)V_0\)
的水，其中
\(0<\alpha<1\)。先向瓶内充入高压气体，高压气体把水在短时间内全部喷出并推动火箭前进。气体视为双原子理想气体，定容摩尔热容
\(C_V=\frac52R\)，大气压为 \(P_0\)。

\begin{enumerate}
\def\labelenumi{\arabic{enumi}.}
\tightlist
\item
  打气机持续提供压强为 \(P_1\)、温度为 \(T_0\)
  的高压气体并缓慢充入塑料瓶。瓶壁很薄且瓶内有水，瓶内气体温度可视为恒定
  \(T_0\)。充到瓶内压强
  \(P_2=\beta P_1\)（\(\beta<1\)）时，求打气机提供的高压气体体积
  \(V_1/V_0\)；
\item
  设单位时间内从打气机打出的体积与压强差满足
  \(\dot V=\gamma\Delta P\)，求充到 \(P_2\) 所需时间；
\item
  发射时气体绝热膨胀并把水在短时间内全部喷出，求喷出水的速度
  \(v_1\)。水密度为 \(\rho_0\)，除水外火箭总质量为 \(M\)；
\item
  取
\end{enumerate}

\[
\alpha=0.75,\quad
\beta=0.20,\quad
\rho_0=1.0\times10^3\,\mathrm{kg/m^3},
\]

\[
M=0.200\,\mathrm{kg},\quad
V_0=2.0\,\mathrm L,\quad
P_1=1.0\times10^6\,\mathrm{Pa},
\]

求第（3）问中水火箭的发射速度 \(v_2\)。

\begin{figure}[H]
\centering
\includegraphics[width=0.678\linewidth,height=\textheight,keepaspectratio,alt={图 7-1　水火箭充气阶段与快速喷射阶段。}]{../figures/question-07.pdf}
\caption*{图 7-1　水火箭充气阶段与快速喷射阶段。}
\end{figure}

\newpage

\section{第 8 题}\label{ux7b2c-8-ux9898}

有一块平凸透镜，凸面为绕 \(y\) 轴旋转的旋转椭球面，轴截面方程为

\[
\frac{x^2}{a^2}+\frac{y^2}{b^2}=1.
\]

用波长为 \(\lambda\) 的单色光沿 \(-y\) 方向入射。

\begin{enumerate}
\def\labelenumi{\arabic{enumi}.}
\tightlist
\item
  把平凸透镜放在一片刚性玻璃板上，在 \(x\ll a,b\)
  范围内得到牛顿环。求从内到外第 \(n\) 条亮条纹的半径 \(r_n\)；
\item
  把透镜放在表面张力很大的透明液体上。液体表面除与透镜接触的部分外近似为圆锥面，液体与透镜完全不浸润。在
  \(x\ll a,b\) 范围内得到牛顿环。求从内到外第 \(n\) 条亮条纹的外半径
  \(r_n'\)。液体表面张力系数为 \(\sigma\)，透镜质量为
  \(m\)，重力加速度为 \(g\)。已知当 \(|z|\ll1\) 时，
\end{enumerate}

\[
(1+z)^\alpha\simeq1+\alpha z.
\]

\begin{figure}[H]
\centering
\includegraphics[width=0.774\linewidth,height=\textheight,keepaspectratio,alt={图 8-1　平凸透镜在刚性玻璃板和透明液面上的放置。}]{../figures/question-08.pdf}
\caption*{图 8-1　平凸透镜在刚性玻璃板和透明液面上的放置。}
\end{figure}
\end{document}
